\documentclass[11pt,a4paper]{article}

\usepackage[utf8]{inputenc}
\usepackage[T1]{fontenc}
\usepackage{mathptmx}
\usepackage[margin=2.5cm]{geometry}
\usepackage{enumitem}
\usepackage[hidelinks]{hyperref}

\linespread{1.15}

\title{Project Report: Personal Portfolio Website}
\author{Ulas Unerli\\ Software Engineering, GISMA University of Applied Sciences}
\date{\today}

\begin{document}

\maketitle
\thispagestyle{empty}

\section{Introduction}
This report documents the design, development, and content of my personal
portfolio website. The portfolio is a single-page website built to present
who I am, the skills I am developing as a software engineering student, and
the ways I can be contacted. The sections below explain where the project can
be found, how the website and its repository are organised, the reasoning
behind the design choices, the tools used to build it, and an honest reflection
on its strengths, weaknesses, and possible future improvements.

\section{Project Links}
The portfolio website and its source code are publicly available at the
following links:

\begin{itemize}
  \item \textbf{Live website:}
    \href{https://ulasunerli.github.io}{https://ulasunerli.github.io}
  \item \textbf{GitHub repository:}
    \href{https://github.com/UlasUnerli/UlasUnerli.github.io}{https://github.com/UlasUnerli/UlasUnerli.github.io}
\end{itemize}

The site is hosted for free using GitHub Pages, which serves the website
directly from the repository, so the repository and the live site always
stay in sync.

\section{Structure of the Website and Repository}

\subsection{Website Structure}
The website is a single HTML page divided into clear sections, with a fixed
navigation menu at the top that links to each one:

\begin{itemize}
  \item \textbf{Header:} displays my name and my role, ``Software Engineer''.
  \item \textbf{Navigation menu:} quick links to About, Skills, Hobbies, and Contact.
  \item \textbf{About:} a short introduction describing who I am and how I work.
  \item \textbf{Skills:} a list of the technologies I use --- HTML, CSS, and Lua.
  \item \textbf{Hobbies:} my interests outside of coding (video games, movies, comics).
  \item \textbf{Contact:} buttons linking to my email, GitHub, and LinkedIn.
  \item \textbf{Footer:} a closing line with my name and the year.
\end{itemize}

\subsection{Repository Structure}
The GitHub repository is intentionally simple. Because the site is built only
with HTML and CSS, the main file is \texttt{index.html}, which contains the page
markup together with the CSS styling inside a \texttt{<style>} block in the
document head. Keeping the styles in the same file keeps the project easy to
manage at this small scale. The repository is named
\texttt{UlasUnerli.github.io}, which is the naming convention GitHub Pages
requires to publish a personal site automatically.

\section{Design Decisions}
The design was guided by one main idea: keep it simple and readable.

\begin{itemize}
  \item \textbf{Single-page layout:} all content fits on one scrollable page so
    a visitor can see everything quickly without navigating between files.
  \item \textbf{Clear navigation:} a sticky menu with anchor links lets the
    reader jump straight to any section.
  \item \textbf{Calm colour scheme:} a dark navy header against a light grey
    background gives good contrast and an easy-to-read, professional look.
  \item \textbf{Sectioned content in cards:} each topic sits in its own white
    box with a border, which separates information clearly and makes the page
    feel organised.
  \item \textbf{Honest, plain content:} the copy reflects my real skill level
    rather than overstating it, which I felt was more genuine and trustworthy.
  \item \textbf{Responsiveness:} a viewport meta tag and a centred, max-width
    layout help the page display reasonably on both desktop and mobile screens.
\end{itemize}

\section{Tools and Technologies}
The portfolio was built using a small, focused set of tools:

\begin{itemize}
  \item \textbf{HTML} --- to structure the content of the page.
  \item \textbf{CSS} --- to style the layout, colours, and spacing.
  \item \textbf{Git \& GitHub} --- for version control and storing the code.
  \item \textbf{GitHub Pages} --- to host the website for free directly from
    the repository.
  \item \textbf{A text/code editor} --- to write and edit the HTML and CSS.
\end{itemize}

I deliberately did not use any frameworks or libraries, since the project is
small and writing plain HTML and CSS gave me a clearer understanding of how
the page works.

\section{Reflection: Strengths, Weaknesses, and Future Work}

\subsection{Strengths}
\begin{itemize}
  \item The site is clean, lightweight, and loads quickly.
  \item The structure is clear and easy to navigate.
  \item The code is simple and readable, which makes it easy to maintain and
    extend.
  \item All contact links work and point to my real profiles.
\end{itemize}

\subsection{Weaknesses}
\begin{itemize}
  \item The skill set shown is currently small, reflecting that I am still
    early in my learning.
  \item There are no live project demonstrations yet, only descriptions.
  \item The visual design is fairly basic and could feel plain to some visitors.
  \item There is no interactivity, since the site uses only HTML and CSS.
\end{itemize}

\subsection{Future Improvements}
\begin{itemize}
  \item Add a dedicated projects section with screenshots and live links as I
    build more.
  \item Learn and introduce JavaScript to add interactive features.
  \item Add a downloadable copy of my CV directly on the page.
  \item Improve the visual design with custom fonts and small animations.
  \item Expand the skills section as I learn new languages and tools.
\end{itemize}

\section{Conclusion}
This portfolio is an honest, simple presentation of my current abilities as a
software engineering student. Building it with plain HTML and CSS helped me
understand the fundamentals of web development and version control, and it gives
me a solid foundation to expand as my skills grow. The reflection above sets a
clear direction for the features I plan to add next.

\end{document}
